\documentclass[10pt,a4paper]{amsart}
\usepackage[margin=19mm]{geometry}
\usepackage{amsmath,amssymb,mathtools}
\usepackage[scr=rsfs,cal=euler]{mathalfa}
\usepackage{xfrac}
\usepackage[unicode,hidelinks,bookmarks=false]{hyperref}
\newcommand{\ie}{i.e.\ }
\newcommand{\cf}{cf.\ }
\newcommand{\resp}{respectively\ }
\newcommand{\Subsection}{\subsection}
\providecommand{\nobreakdash}{\nobreak\hbox{-}}
\setlength{\emergencystretch}{2em}
\multlinegap0pt
\newtheorem{problem}{Problem}
\newtheorem{theorem}{Theorem}
\newcommand{\E}{\mathcal{E}}
\newcommand{\R}{\mathbb{R}}
\newcommand{\U}{\mathcal{U}}
\newcommand{\X}{\mathcal{X}}
\title{\LARGE \bf Exploiting Over-Approximation Errors as Preview Information for Nonlinear Control*}
\author{Antoine Aspeel, Antoine Girard, Thiago Alves Lima}
\date{Source: arXiv:2511.03577v2 (CC BY 4.0)}
\begin{document}
\maketitle

\noindent\emph{Source and licence.} \LARGE \bf Exploiting Over-Approximation Errors as Preview Information for Nonlinear Control*. Version arXiv:2511.03577v2 (CC BY 4.0), distributed by arXiv under the Creative Commons Attribution 4.0 International licence, http://creativecommons.org/licenses/by/4.0/, as recorded in the arXiv licence metadata for this exact version and shown on the abstract page https://arxiv.org/abs/2511.03577v2. Full licence notice: \texttt{licenses/arxiv-2511.03577-CC-BY-4.0.txt}.\par\smallskip

\noindent\textit{Editorial scope.} Counted unit: the article's theorem thm:overapproximation (source0.tex lines 166--171). The article's complete original proof of it is delivered with it (source0.tex lines 172--180, one \texttt{proof} environment of 9 lines). No mathematics is written, repaired or re-derived: the statement and the proof are the article's own lines, in the article's own notation, and no retained line is substituted or re-flowed.  Retained uncounted as the source-local context the proof needs: lines 103--105; lines 127--132.  Nothing was omitted from any retained span, so the package carries no omission record.  No retained line contains a citation, so the wrapper carries no bibliography and no bibliography entry.  No retained line contains a figure environment, an image-inclusion command or drawing code, so no image is delivered and the package declares no asset dependency.  Editorial formatting only, in addition: an AMS \texttt{amsart} preamble that restates the article's own declarations and macros used by the retained lines, namely the environments \texttt{problem}, \texttt{theorem} on separate counters, together with the standard packages those lines need.  The retained environments print in this document as Problem 1, Theorem 1 in order of appearance, whereas the article prints the same items with the numbering of its own full document.  The complete native source of the article as distributed in the arXiv e-print for this exact version is bundled unchanged as \texttt{sources/188520/source0.tex}.
\begin{align}\label{eq:errorInclusion}
\text{ for all } (x,u)\in\X\times\U:\ e(x,u)\in\E,
\end{align}

\begin{problem}[Concretization]\label{prob:concretization}
Find $u\in\U$ such that
$$
u=\pi(x,f(x,u)-\hat{f}(x,u)).
$$
\end{problem}

\begin{theorem} \label{thm:overapproximation}
Let $f,\hat{f}:\X\times\U\rightarrow\R^{n_x}$ and $\pi:\X\times\E\rightarrow\U$ be three functions, and assume that $\E$ satisfies \eqref{eq:errorInclusion}. For any $x\in\X$, if Problem~\ref{prob:concretization} has a solution $u$, then
$$
f(x,u)\in \{ F(x,\pi(x,\bar{e}),\bar{e}) \mid \bar{e}\in\E \}.
$$
\end{theorem}

\begin{proof}
Consider $\bar{e} \coloneq f(x,u)-\hat{f}(x,u)$ and note that $\bar{e}\in\E$ thanks to \eqref{eq:errorInclusion}. In addition, since $u$ is a solution to Problem~\ref{prob:concretization}, $u=\pi(x,\bar{e})$. Finally, one can write
\begin{align*}
f(x,u) &= \hat{f}(x,u)+f(x,u)-\hat{f}(x,u) \\
&=\hat{f}(x,\pi(x,\bar{e}))+\bar{e} \\
&= F(x,\pi(x,\bar{e}),\bar{e}),
\end{align*}
which concludes the proof.
\end{proof}

\end{document}
